\begin{lemma}[Localized weak-type $(1,1)$ bound for the maximal truncation]
Fix $n \geq 1$ and $j \in \{1,\ldots,n\}$, let $\lambda$ be Lebesgue measure on
$\mathbb{R}^n$, let $\gamma_n$ be the Gaussian measure of
\noderef{GaussMeasure}, let $K_{\gamma,j}$ be the kernel of
\noderef{GaussRieszKernel}, and let
\[
  T : L^1(\gamma_n;\mathbb{R}) \longrightarrow
      \mathcal{M}\bigl(\gamma_n;\mathbb{R}^n\bigr)
\]
be \emph{any} linear map into the $\gamma_n$-almost-everywhere equivalence
classes of measurable $\mathbb{R}^n$-valued functions such that
\begin{enumerate}
\item[(T1)] $T(\iota v) = R_\gamma v$ for every $v \in L^2(\gamma_n)$, with
$\iota$ the inclusion of \noderef{ToL1} and $R_\gamma$ the transform of
\noderef{RieszTransform}; and
\item[(T2)] $t\,\gamma_n\bigl(\{x : |Tf(x)|_{\mathbb{R}^n} > t\}\bigr)
\leq 2\,\|f\|_{L^1(\gamma_n)}$ for every $f \in L^1(\gamma_n;\mathbb{R})$ and
every real $t > 0$.
\end{enumerate}
For $f \in L^1(\gamma_n)$, written also for its canonical representative, put
\[
  T_{j,\varepsilon}f(x) = \int_{|x-y|>\varepsilon}K_{\gamma,j}(x,y)f(y)\,dy ,
  \qquad
  T_{j,*}f(x) = \sup_{\varepsilon>0}\bigl|T_{j,\varepsilon}f(x)\bigr|
  \;\in\;[0,+\infty] ,
\]
the $y$-integral against Lebesgue measure and the supremum taken in the
extended nonnegative reals.  Let $R \geq 1$.  Then there is a constant
$C \in (0,\infty)$, depending only on $n$, $j$ and $R$, such that for every
$f \in L^1(\gamma_n)$ and every $t \in (0,+\infty]$,
\[
  \gamma_n\Bigl(\bigl\{\,x \in \mathbb{R}^n :\ |x| \leq R
     \ \text{ and }\ T_{j,*}f(x) > t\,\bigr\}\Bigr)
  \;\leq\; \frac{C}{t}\,\|f\|_{L^1(\gamma_n)} ,
\]
the right-hand side read in $[0,+\infty]$.

\emph{This is the reference paper's displayed maximal estimate, localized.}  The
paper states it with no restriction on $x$ and with a constant $C_{n,j}$
depending only on the dimension and the coordinate; the present node restricts
$x$ to the closed ball of radius $R$ and allows the constant to depend on $R$ as
well.  That is the whole of the departure, and it is recorded, with its
return-to-faithful argument, in the deviation file
\texttt{reference/maximal-truncation-localization.tex}, which this node claims.
Nothing else about the statement differs: the operator, the truncations, the
exponent and the shape $C\|f\|_1/t$ are the paper's.

\emph{The supremum and the threshold are extended-real, and this is
load-bearing.}  $T_{j,*}f(x)$ is a supremum over $\varepsilon > 0$ of real
numbers which is not known in advance to be finite --- the far-majorant bound
$|T_{j,\varepsilon}f(x)| \leq C(\varepsilon)\|f\|_{L^1(\gamma_n)}$ degrades as
$\varepsilon \downarrow 0$ --- so it is taken in $[0,+\infty]$, exactly as in
\noderef{GaussRieszCotlarLocal}, whose left-hand side this is verbatim.  A real
supremum would take a conventional value where the family is unbounded, which
is precisely where the estimate carries content.  Correspondingly the threshold
$t$ ranges over $(0,+\infty]$: at $t = +\infty$ the level set is empty and the
assertion is the triviality $0 \leq 0$, and the content is at finite $t$.

\emph{No measurability of the level set is asserted or needed.}  $\gamma_n$ is
applied to the level set as an outer measure, which is defined on every subset
and is monotone and countably subadditive there.  That is enough for the proof
below, because every inequality it chains together --- the weak-type bounds for
$M_2$ and the comparison $\gamma_n \leq \pi^{-n/2}\lambda$ --- is itself stated
for arbitrary sets.  Bounding the outer measure is the stronger assertion, so a
consumer that does know the set to be measurable loses nothing.

\emph{The constant depends on $R$, and the object is coordinatewise.}  $C$
degrades in $R$ through the Cotlar constant and through
$a_{R+2} = \pi^{-n/2}e^{-(R+2)^2}$, the lower bound for the Gaussian density on
$\overline B(0,R+2)$; no uniformity in $n$, $j$ or $R$ is claimed, and the paper
explicitly licenses the dimension dependence at this step.  The quantity
bounded is the coordinatewise maximal truncation
$\sup_\varepsilon|T_{j,\varepsilon}f|$, not the paper's vector maximal
truncation $\sup_\varepsilon|T_\varepsilon f|_{\mathbb{R}^n}$; nothing here
asserts anything about the latter.

\emph{The hypothetical style.}  (T1) and (T2) are written out in full and not
imported by reference; the lemma is a statement about every $T$ satisfying
them, as \noderef{CoordinateEndpointBound} and \noderef{GaussRieszCotlarLocal}
are, and it is not vacuous because \noderef{L1Extension} exhibits such a $T$.
\end{lemma}
\begin{proof}
Write $p = \pi^{-n/2}$, $a = \pi^{-n/2}e^{-(R+2)^2} \in (0,\infty)$ --- so that
$a^{-1} = \pi^{n/2}e^{(R+2)^2}$ --- and $E = e^{(R+2)^2}$.  Throughout,
$\overline B = \overline B(0,R+2)$.

\emph{Step 0: the Cotlar constant.}
Apply \noderef{GaussRieszCotlarLocal} to this $T$ (with the two properties it
requires, which are (T1) and (T2) and hold by assumption), to $j$, and to this
$R$, whose hypothesis $R \geq 1$ holds.  It supplies a constant
$C_0 \in (0,\infty)$ --- in particular $C_0 \neq 0$ and $C_0 \neq +\infty$ ---
such that for every $f \in L^1(\gamma_n)$ and every $x$ with $|x| \leq R$,
\begin{equation*}
  T_{j,*}f(x) \;\leq\; C_0\bigl[\,\Psi_f(x)^2 + \Phi_f(x) + N_f\,\bigr] ,
  \tag{1}
\end{equation*}
where, with $M_2$ the local maximal function of \noderef{LocalMaximalFunction},
\[
  \Phi_f = M_2\bigl(w_f\bigr),\quad w_f = f\mathbf 1_{\overline B};
  \qquad
  \Psi_f = M_2\bigl(u_f\bigr),\quad u_f = |(Tf)_j|^{1/2}\mathbf 1_{\overline B};
  \qquad
  N_f = \|f\|_{L^1(\gamma_n)} .
\]
Crucially $C_0$ does not depend on $f$: it is produced before $f$ is quantified.
Define
\[
  C = 3\,C_0\bigl(1 + 33\cdot4^n E\bigr) .
\]
Then $C \in (0,\infty)$ and $C$ depends only on $n$, $j$ and $R$.  We show the
lemma with this $C$.

Fix $f \in L^1(\gamma_n)$ and $t \in (0,+\infty]$, and abbreviate
$\Phi = \Phi_f$, $\Psi = \Psi_f$, $N = N_f$, $w = w_f$, $u = u_f$.  Note
$N < \infty$.  Put
\[
  s = \frac{t}{3C_0} ,
\]
which is nonzero because $t \neq 0$ and $3C_0 < \infty$.

\emph{Step 1: measurability of the two inputs.}
The canonical representative of an almost-everywhere class is measurable, so
$(Tf)_j$ is measurable; hence $u = |(Tf)_j|^{1/2}\mathbf 1_{\overline B}$ is
measurable, the square root being continuous on $[0,\infty)$ and
$\overline B$ closed.  For $w$: the canonical representative $f$ is
$\gamma_n$-almost everywhere equal to a measurable function, and by clause (4)
of \noderef{GaussLebesgueBallComparison} the measures $\gamma_n$ and $\lambda$
have the same null sets, so $f$ is $\lambda$-almost everywhere equal to a
measurable function; multiplying by the indicator of the measurable set
$\overline B$ preserves this.  Thus $w$ and $u$ are Lebesgue-almost-everywhere
measurable, which is the hypothesis both weak-type nodes require.

\emph{Step 2: the level set is covered by three sets.}
Let
\[
  \mathcal{E} = \bigl\{x : |x| \leq R \text{ and } T_{j,*}f(x) > t\bigr\},
  \qquad
  S_1 = \{x : \Phi(x) > s\},\quad
  S_2 = \{x : \Psi(x)^2 > s\},\quad
  S_3 = \{x : N > s\} .
\]
We claim $\mathcal{E} \subseteq S_1 \cup S_2 \cup S_3$.  Indeed, let
$x \in \mathcal{E}$ and suppose $x$ lay in none of the three.  Then
$\Psi(x)^2 \leq s$, $\Phi(x) \leq s$ and $N \leq s$, so
$\Psi(x)^2 + \Phi(x) + N \leq 3s = 3\,t/(3C_0) = t/C_0$, the cancellation of
the factor $3$ being legitimate in $[0,+\infty]$ because $3 \in (0,\infty)$.
Multiplying by $C_0 \in (0,\infty)$ and using $(1)$,
$T_{j,*}f(x) \leq C_0\,(t/C_0) = t$, contradicting $T_{j,*}f(x) > t$.  Hence
by monotonicity and subadditivity of the outer measure $\gamma_n$ (shared
context),
\begin{equation*}
  \gamma_n(\mathcal{E}) \leq \gamma_n(S_1) + \gamma_n(S_2) + \gamma_n(S_3) .
  \tag{2}
\end{equation*}
Note that $S_3$ does not depend on $x$: it is $\mathbb{R}^n$ if $N > s$ and
$\emptyset$ otherwise.

\emph{Step 3: the term $S_1$.}
By \noderef{LocalMaximalWeakType}, applied to $w$ (measurable up to a null set
by Step 1) at the level $s \neq 0$,
\[
  \lambda(S_1) \leq \frac{4^n}{s}\int_{\mathbb{R}^n}|w|\,d\lambda .
\]
By \noderef{GaussMeasure}, $\gamma_n$ is $\lambda$ weighted by the density
$\varrho(y) = \pi^{-n/2}e^{-|y|^2}$, and $\varrho \geq a$ on $\overline B$
because $|y| \leq R+2$ there; since $w$ vanishes off $\overline B$,
\[
  a\int_{\mathbb{R}^n}|w|\,d\lambda
    = a\int_{\overline B}|f|\,d\lambda
    \leq \int_{\overline B}|f|\,\varrho\,d\lambda
    = \int_{\overline B}|f|\,d\gamma_n \leq N ,
\]
so $\int|w|\,d\lambda \leq a^{-1}N = \pi^{n/2}E\,N$.  Finally, by clause (1) of
\noderef{GaussLebesgueBallComparison}, which holds for \emph{every} set,
$\gamma_n(S_1) \leq p\,\lambda(S_1)$.  Combining, and using
$p\,\pi^{n/2} = 1$,
\begin{equation*}
  \gamma_n(S_1) \leq \frac{p\,4^n\pi^{n/2}E\,N}{s}
    = \frac{4^nE\,N}{s} = \frac{3C_0\,4^nE\,N}{t} . \tag{3}
\end{equation*}

\emph{Step 4: the term $S_2$.}
Let $\sqrt s \in (0,+\infty]$ be the square root of $s$, i.e.\ the unique
element of $[0,+\infty]$ whose square is $s$; it is nonzero because $s$ is.
Squaring is strictly monotone on $[0,+\infty]$, so $\Psi(x)^2 > s = (\sqrt s)^2$
forces $\Psi(x) > \sqrt s$, whence
$S_2 \subseteq \{x : \Psi(x) > \sqrt s\}$.  By
\noderef{LocalMaximalWeakL2}, applied to $u$ (measurable by Step 1) at the
level $\sqrt s \neq 0$, and using $(\sqrt s)^2 = s$,
\begin{equation*}
  \lambda(S_2) \leq \frac{16\cdot4^n}{s}\,\mathcal{A},
  \qquad
  \mathcal{A} = \sup_{0<\sigma\leq+\infty}
      \sigma^2\,\lambda\bigl(\{y : |u(y)| > \sigma\}\bigr) . \tag{4}
\end{equation*}

We claim the \emph{uniform} bound
\begin{equation*}
  \mathcal{A} \leq 2\,a^{-1}N = 2\,\pi^{n/2}E\,N . \tag{5}
\end{equation*}
It suffices to bound each term of the supremum by the right-hand side.  If
$\sigma = +\infty$ the set $\{|u| > \sigma\}$ is empty, $u$ being real valued,
and the term is $(+\infty)\cdot0 = 0$.  So let $\sigma$ be a real number with
$\sigma > 0$ and put
\[
  E_\sigma = \{y : |u(y)| > \sigma\}
    = \{y \in \overline B : |(Tf)_j(y)| > \sigma^2\} ,
\]
the second description because $u = |(Tf)_j|^{1/2}$ on $\overline B$ and $0$
off it, and because $r \mapsto r^{1/2}$ is strictly increasing on $[0,\infty)$.
The set $E_\sigma$ is measurable (Step 1) and contained in
$\overline B = \overline B(0,R+2)$, so clause (2) of
\noderef{GaussLebesgueBallComparison}, \emph{instantiated at $R+2$ in place of
its $R$}, gives $a\,\lambda(E_\sigma) \leq \gamma_n(E_\sigma)$, hence
$\lambda(E_\sigma) \leq a^{-1}\gamma_n(E_\sigma)$ since $a \in (0,\infty)$.
Next, $E_\sigma \subseteq \{y : |(Tf)_j(y)| > \sigma^2\}$, and
\noderef{CoordinateEndpointBound} --- applied to this $T$ with its two
properties, to $f$, to the coordinate $j$ and to the real level
$\sigma^2 > 0$ --- gives
\[
  \sigma^2\,\gamma_n\bigl(\{y : |(Tf)_j(y)| > \sigma^2\}\bigr) \leq 2N .
\]
Therefore
\[
  \sigma^2\lambda(E_\sigma)
    \leq a^{-1}\,\sigma^2\gamma_n(E_\sigma)
    \leq a^{-1}\,\sigma^2\gamma_n\bigl(\{|(Tf)_j| > \sigma^2\}\bigr)
    \leq 2\,a^{-1}N ,
\]
a bound independent of $\sigma$, which is $(5)$.  (The uniformity is the point:
a bound still depending on $\sigma$ would not bound the supremum.)

Feeding $(5)$ into $(4)$ and applying clause (1) of
\noderef{GaussLebesgueBallComparison} again,
\begin{equation*}
  \gamma_n(S_2) \leq p\,\lambda(S_2)
    \leq \frac{p\cdot16\cdot4^n\cdot2\pi^{n/2}E\,N}{s}
    = \frac{32\cdot4^nE\,N}{s}
    = \frac{96\,C_0\,4^nE\,N}{t} . \tag{6}
\end{equation*}

\emph{Step 5: the term $S_3$, and the conclusion.}
If $N \leq s$ then $S_3 = \emptyset$ and $\gamma_n(S_3) = 0$.  Otherwise
$N > s$; then $s \neq +\infty$ (nothing exceeds $+\infty$), so
$s \in (0,\infty)$ and $N/s > 1$, while $\gamma_n(S_3) = \gamma_n(\mathbb{R}^n)
 = 1$ because $\gamma_n$ is a probability measure (\noderef{GaussIsProb}).
In both cases
\begin{equation*}
  \gamma_n(S_3) \leq \frac{N}{s} = \frac{3C_0\,N}{t} . \tag{7}
\end{equation*}

Adding $(3)$, $(6)$ and $(7)$ into $(2)$,
\[
  \gamma_n(\mathcal{E})
    \leq \bigl(3C_0\,4^nE + 96\,C_0\,4^nE + 3C_0\bigr)\frac{N}{t}
    = 3C_0\bigl(1 + 33\cdot4^nE\bigr)\frac{N}{t}
    = \frac{C\,N}{t} ,
\]
which is the assertion.
%
% FORMALIZATION NOTE (cut points for ProofFormalization, recorded while the
% argument is fresh; decomposition is free once a child's Lean closes, and is
% not to be done in the stating phase).  Two natural children:
%   (D1) Step 4's uniform bound (5): for every sigma in (0, +infinity],
%        sigma^2 lambda({|u| > sigma}) <= 2 a^{-1} ||f||_1, where
%        u = |(Tf)_j|^{1/2} 1_{closedBall 0 (R+2)}.  This is the only place
%        CoordinateEndpointBound and GaussLebesgueBallComparison clause (2)
%        appear, and it is the step the reviewer flagged as delicate.  Do NOT
%        make it a tablet node in a stating phase: as a standalone statement it
%        is close to a repackaging of CoordinateEndpointBound composed with the
%        ball comparison, so it is exposed to Substantiveness Clause 2.
%   (D2) Step 2's three-way split: from u <= C(A + B + c) and t < u conclude
%        A > t/(3C) or B > t/(3C) or c > t/(3C), in ENNReal with C in (0,inf).
%        Pure ENNReal arithmetic, reusable, and the place an informal proof is
%        most likely to slip.
% Steps 3 and 5 are each four lines and need no cut.
\end{proof}
